\documentclass[11pt]{template}

% Optional direkt außerhalb des Templates:
% \documentclass[10pt]{template}
% \documentclass[12pt]{template}
%
% Optional: Abstand zwischen Caption und Inhalt anpassen:
% \setcaptiongap{1pt}

\title{Praktikum Grundlagen der 3D Druck Technologien}
\subtitle{Bericht}
\author{Nils Würtz}
\matrikelnummer{6376725}
\fachbereich{3}
\studiengang{Informatik Vollfach}
\semester{Sommersemester 2026}

\begin{document}

\maketitle
\makefrontmatter

\chapter{Projektbeschreibung}

\section{Projektumfeld}
Hier beginnt der eigentliche Text.
Hierbei ist es gutasdasdasd ada dasd adbadvakd vad ald adva sd adad odosdosobd sobd ob sob

asd

\section{Projektziel}
Eine \ac{api} wird automatisch ausgeschrieben.

\subsection{Abgrenzung}
Beispieltext mit einer Fußnote.\footnote{Diese Fußnote wird mit 9 pt gesetzt.}

\chapter{Projektvorbereitung}

\section{Ist-Analyse}

\begin{figure}[H]
  \centering
  \caption{Beispielabbildung}
  \fbox{\rule{0pt}{3cm}\rule{5cm}{0pt}}
  \source{Eigene Darstellung}
  \label{fig:example}
\end{figure}

\section{Soll-Analyse}

\begin{table}[H]
  \centering
  \caption{Beispieltabelle}
  \begin{tabular}{ll}
    \textbf{Eigenschaft} & \textbf{Wert} \\
    Beispiel             & 42
  \end{tabular}
  \source{Eigene Darstellung}
  \label{tab:example}
\end{table}


\section{Quellcode}

\begin{documentcode}{CSharp}{Initialisierung eines Services}{Eigene Darstellung}
public sealed class AppSdkService
{
    public void Initialize()
    {
        Console.WriteLine("Initialized");
    }
}
\end{documentcode}

\begin{documentcode}{TypeScript}{Beispiel einer TypeScript-Funktion}{Eigene Darstellung}
function add(a: number, b: number): number {
    return a + b;
}
\end{documentcode}

% Weitere mögliche Sprachen:
% Java, C, C++, Python, SQL, XML, JavaScript, JSON, CSharp, TypeScript

% Harvard:
% parenthetisch: (Knuth, 1984)
Eine Beispielquelle ist \parencite{knuth1984}.

% Mit Seitenangabe: (Knuth, 1984, S. 12)
Ein konkreter Seitenverweis ist \parencite[12]{knuth1984}.

% Textuell: Knuth (1984)
\textcite{knuth1984} beschreibt ...

\printreferences

\startappendices

\appendixentry{Quellcode XYZ}
Hier steht der Quellcode oder ein Verweis darauf.

\appendixentry{Physischer Netzwerkplan}
Hier steht der Netzwerkplan.

\end{document}
